% page 354 du fac-simile (image p-353 du scan)
% bloc 154.9 x 229.0 mm ; marge gauche 8.0 mm
\pgc{-12.700mm}{0.000mm}{
\l{\cel{3}346}
\l{}
\l{}
\l{\sou{luktar}. (\sou{netrans., kontre, kun, por}).I. (\sou{aludante du}\cel{1}in-}
\l{\cel{3}\sou{dividui}). Esforcar por renversar la altro, per em-}
\l{\cel{3}braco. - II. (\sou{aludante du individui, du populi,}\cel{1}du}
\l{\cel{3}\sou{partisi}). Esforcar por vinkar la altra, la adverso,}
\l{\cel{3}kontestante pri sua superioreso. - III. (\sou{aludante}\cel{1}la}
\l{\cel{3}\sou{volo da la homo}).Esforcar por vinkar sua pasioni, por}
\l{\cel{3}vinkar la obstakli extera.- IV. (\sou{aludante forco}).}
\l{\cel{3}Efikar ad altra forco, kontre-sinse. - FIS.}
\l{}
\l{\sou{lular}. (\sou{trans.)}\cel{2}Duktar (infanteto) ad dormo per kantar apud}
\l{\cel{3}lu.\cel{2}DE.}
\l{}
\l{\sou{lum}ar. (\sou{netrans}.)(\sou{aludante ula korpi}).Emisar ula radii qui}
\l{\cel{3}igas videbla la objekti. - EFIS.}
\l{}
\l{\sou{lumbo}. (\sou{anat.}) Dopa parto di la torso, de la dorso til la}
\l{\cel{3}hanchi, dextre e sinistre di la kolono vertebrala (spino).}
\l{\cel{3}- EFIS.}
\l{}
\l{\sou{lumbago}. (patol.) Doloro en la regiono lumbala, rumatismala,}
\l{\cel{3}nevralgiala, o produktita da esforco tro granda. -EFIS.}
\l{}
\l{\sou{lumpo}. (\sou{mineral.}) Fragmento di roko ronda quan onu ren-}
\l{\cel{3}kontras en ula strati mineralala. - E.}
\l{}
\l{\sou{luno}. (\sou{astron.)}\cel{2}Planeto qua jiras cirkum la Terglobo}
\l{\cel{3}(di qua lu nomizesis la \textquotedbl{}satelito\textquotedbl{}) e qua lumizas lu}
\l{\cel{3}dum nokto kande lu esas situita tale ke lu reflektas}
\l{\cel{3}a la Terglobo la lumo-radii da Suno. - DEFIRS.}
\l{}
\l{\sou{lunatika}. Submisita a la efiko (asertita) da Luno ; dic-}
\l{\cel{3}esas precipue pri persono qua ne plus havas sua ra-}
\l{\cel{3}ciono o rezono-fakultato pro, asertite, subisir la}
\l{\cel{3}efiko da Luno. - DEFIRS.}
\l{}
\l{\sou{lun}dio. La dio duesma di la semano (olta qua sequas}
\l{\cel{3}sundio). - FIS.}
\l{}
\l{\sou{lupo}. (\sou{optiko)}. Lenso bikonvexa, uzata por igar aspektar}
\l{\cel{3}plu granda kam reale la objekto quan onu regardas tra}
\l{\cel{3}ta lenso. - DF.}
\l{}
\l{\sou{lupanaro}.\cel{2}Domo en qua la mulieri su prostitucas. - EFL.}
\l{}
\l{\sou{lupio}. (patol.)Surkreskajo sub la pelo, tumoro sen-dolora,}
\l{\cel{3}maxim-multa-kaze kisteskinta. - FIS.}
\l{}
\l{\sou{lupino.}\cel{2}(\sou{bot.)}\cel{2}Planto leguminosa, di qua la shelo kon-}
\l{\cel{3}tenas grani di qui la farino esas manjebla.- L. lu-}
\l{\cel{3}\sou{pinus}. - DEFIRS.}
}
